\itemhead{第2章 小结与复习 · 第3题（教材p.60）}{XJ-B1-C02-SU-P060-Y1169}
\label{XJ-B1-C02-SU-P060-Y1169}
{\footnotesize 湘教版教材p.60 / PDF第67页 / 第2章 小结与复习 / 第3题}
\field{原题}{基本不等式以及不等式的性质是证明不等式的重要手段和方法。你能证明基本不等式吗？在使用基本不等式解决最大（小）值等实际问题时，需要注意什么？}
\field{参考答案}{由 $(\sqrt a-\sqrt b)^2\ge0$ 得 $a+b\ge2\sqrt{ab}$。求最值须检查正数条件、和或积是否固定、等号条件能否在定义域内同时实现。}
{\small \field{关联知识}{I03 基本不等式、等号与最值建模（教材帮p.52、53、54 / PDF61、62、63）}
\field{方法与易错点}{列正数条件与固定量，构造基本不等式，解等号条件并代回实际限制。 易错：把一个下界误当最小值；多个等号条件不能同时实现；漏单位和整数限制。}
\field{答案依据}{答案册未提供本栏目的独立对应答案；依据原题、教材概念补写参考回应。开放性总结不设唯一数值答案。}}
